\documentclass[11pt]{article}
\usepackage{amsmath,amssymb,amsthm}
\usepackage{hyperref}
\usepackage[margin=1in]{geometry}

\newcommand{\N}{\mathbb{N}}
\newcommand{\Z}{\mathbb{Z}}
\newcommand{\Q}{\mathbb{Q}}
\newcommand{\La}{\Lambda}
\newcommand{\GK}{\operatorname{GK}}
\newcommand{\CPP}{\operatorname{CPP}}
\newcommand{\ALCD}{\operatorname{ALCD}}
\newcommand{\IP}{\operatorname{IP}}
\newcommand{\qpoch}[2]{(#1;#2)_\infty}
\newcommand{\qpochn}[3]{(#1;#2)_{#3}}
\newcommand{\qbin}[2]{\genfrac{[}{]}{0pt}{}{#1}{#2}}
\DeclareMathOperator{\tr}{tr}

\theoremstyle{plain}
\newtheorem{theorem}{Theorem}
\newtheorem{conjecture}[theorem]{Conjecture}
\newtheorem{proposition}[theorem]{Proposition}
\newtheorem{lemma}[theorem]{Lemma}
\newtheorem{corollary}[theorem]{Corollary}
\theoremstyle{definition}
\newtheorem{definition}[theorem]{Definition}
\newtheorem{example}[theorem]{Example}
\newtheorem{remark}[theorem]{Remark}

\title{The $Q$-Operator and its Adjoint:\\
From Level~3 Operator Traces to Cascade Commutators}
\author{Robin Langer}
\date{\today}

\begin{document}
\maketitle

\begin{abstract}
We develop two complementary operator-theoretic frameworks for the
cylindric Kostka polynomials $Q_{n,c}(q)$.

\smallskip\noindent
\textbf{Part~I} lifts the construction to Level~3 in the hierarchy
of generating functions: the operator trace framework with free
parameters.  This produces a genuine operator
$\hat{Q}_{c,\mathbf{n}}$ on the ring of symmetric functions~$\La$
whose trace recovers the multivariate cylindric Kostka polynomial.
The adjoint $\hat{Q}_{c,\mathbf{n}}^*$ corresponds to the
reversed-complementary profile, and we show that
$\tr(\hat{Q}^*) = \tr(\hat{Q})$: the adjoint is invisible at
the polynomial level.

\smallskip\noindent
\textbf{Part~II} defines a second operator $\mathbf{Q}(z)
= \sum Q_m z^m P_m$ that packages the $Q$-polynomials as
coefficients of Pieri multiplication operators.  Computing the
products $\mathbf{Q}(z)\mathbf{Q}^*(w)$ and
$\mathbf{Q}^*(w)\mathbf{Q}(z)$ reveals cross-correlation
functions and a cascade commutator structure that couples
different depth levels of the cylindric partition.
\end{abstract}

\tableofcontents


%% ============================================================
%% PART I: THE LEVEL 3 FRAMEWORK
%% ============================================================

\part{The Level~3 Operator Trace Framework}


%% ============================================================
\section{The Level~3 operator}
\label{sec:level3}
%% ============================================================

\subsection{Setup}

Let $\La$ denote the ring of symmetric functions with Hall inner
product $\langle\,\cdot\,,\,\cdot\,\rangle$ satisfying
$\langle s_\lambda, s_\mu \rangle = \delta_{\lambda,\mu}$,
or equivalently $\langle h_\lambda, m_\mu \rangle
= \delta_{\lambda,\mu}$.  The Cauchy identity reads:
\begin{equation}\label{eq:cauchy}
  \sum_\lambda s_\lambda(X)\, s_\lambda(Y)
  = \prod_{i,j} \frac{1}{1 - x_i y_j}
  = \Omega[XY].
\end{equation}

For rank~3 with profile $c = (c_0, c_1, c_2)$, set
$T := d = c_0 + c_1 + c_2$, and let
$\pi = (\pi_0, \ldots, \pi_{T-1}) \in \{0,1\}^T$ be the binary
string encoding the profile (with $n$~ones and $m$~zeros,
$n + m = T$).

\subsection{The vertex operators}

Define operators on $\La$ depending on a formal parameter~$u$:
\begin{align}
  G^1(u) &:= \Omega[Xu]
    \quad\text{(Pieri addition: $s_\mu \mapsto \sum_{\lambda \in
    U(\mu)} s_\lambda\, u^{|\lambda|-|\mu|}$)},
    \label{eq:G1} \\
  G^0(u) &:= \Omega^*[Xu]
    \quad\text{(Pieri removal: the Hall-adjoint of $G^1(u)$)}.
    \label{eq:G0}
\end{align}
In terms of the translation and multiplication operators of
\cite{COMMUTATION}:
\begin{equation}\label{eq:G-PT}
  G^1(u) = P(u) = \mathcal{P}(u), \qquad
  G^0(u) = T(u) = \mathcal{T}(u) = P^*(u).
\end{equation}

The \emph{degree operator} $D_w$ acts diagonally in the Schur basis:
\begin{equation}\label{eq:Dw}
  D_w\, s_\lambda = w^{|\lambda|}\, s_\lambda.
\end{equation}

\subsection{The operator product}

\begin{definition}[Level~3 operator]
\label{def:O-hat}
The \emph{Level~3 operator} for profile $\pi$ with parameters
$\mathbf{u} = (u_0, \ldots, u_{T-1})$ and $w$ is:
\begin{equation}\label{eq:O-hat}
  \hat{O}_c(\mathbf{u}, w)
  := G^{\pi_0}(u_0)\, G^{\pi_1}(u_1) \cdots
     G^{\pi_{T-1}}(u_{T-1})\, D_w
  \;\colon\; \La \to \La.
\end{equation}
\end{definition}

The multivariate generating function for cylindric plane partitions
is recovered as the trace:
\begin{equation}\label{eq:trace}
  \GK_c(\mathbf{z})
  = \tr(\hat{O}_c)
  := \sum_\mu \langle s_\mu,\, \hat{O}_c\, s_\mu \rangle,
\end{equation}
under the specialisation
\begin{equation}\label{eq:specialisation}
  w = z_0 z_1 z_2, \qquad
  u_k = \begin{cases}
    z_0 z_1 \cdots z_{k-1} & \text{if } \pi_k = 1, \\
    (z_0 z_1 \cdots z_{k-1})^{-1} & \text{if } \pi_k = 0.
  \end{cases}
\end{equation}

\begin{remark}[Matrix elements]
\label{rem:matrix-elements}
The trace uses only the \emph{diagonal} matrix elements
$\langle s_\mu, \hat{O}_c\, s_\mu \rangle$.
The \emph{off-diagonal} elements
$\langle s_\lambda, \hat{O}_c\, s_\mu \rangle$ with
$\lambda \neq \mu$ count ``open'' interlacing sequences from~$\mu$
to~$\lambda$ (without the cylindric periodicity condition
$\mu^0 = \mu^T$).  The trace enforces periodicity.
\end{remark}


%% ============================================================
\section{The adjoint operator}
\label{sec:adjoint}
%% ============================================================

\subsection{Adjoint of the building blocks}

With respect to the Hall inner product, the vertex operators are
mutual adjoints:
\begin{equation}\label{eq:G-adjoint}
  (G^1(u))^* = G^0(u), \qquad (G^0(u))^* = G^1(u).
\end{equation}
This is the operator-theoretic content of the Pieri rules: adding a
horizontal strip is adjoint to removing one.

The degree operator is self-adjoint:
\begin{equation}\label{eq:Dw-adjoint}
  D_w^* = D_w,
\end{equation}
since $\langle D_w s_\lambda, s_\mu \rangle
= w^{|\lambda|} \delta_{\lambda\mu}
= \langle s_\lambda, D_w s_\mu \rangle$.

\subsection{Adjoint of $\hat{O}_c$}

Taking the adjoint of the product \eqref{eq:O-hat} reverses the
order and replaces each factor by its adjoint:
\begin{equation}\label{eq:O-adjoint}
  \hat{O}_c^*
  = D_w\; G^{\bar\pi_{T-1}}(u_{T-1})\;
    G^{\bar\pi_{T-2}}(u_{T-2}) \cdots G^{\bar\pi_0}(u_0),
\end{equation}
where $\bar\pi_k := 1 - \pi_k$ is the complementary bit:
addition~$\leftrightarrow$ removal.

\begin{remark}[Profile transformation]
\label{rem:profile-transform}
The adjoint $\hat{O}_c^*$ is the Level~3 operator for the
\emph{reversed-complementary profile}
$\bar\pi^{\mathrm{rev}} = (\bar\pi_{T-1}, \bar\pi_{T-2}, \ldots,
\bar\pi_0)$ with reversed parameters
$\mathbf{u}^{\mathrm{rev}} = (u_{T-1}, \ldots, u_0)$, up to the
position of~$D_w$.
\end{remark}


\subsection{Behaviour under the trace}

\begin{proposition}[Trace invariance]
\label{prop:trace-invariance}
$\tr(\hat{O}_c^*) = \tr(\hat{O}_c)$.
\end{proposition}

\begin{proof}
Since the Hall inner product is symmetric
($\langle f, g \rangle = \langle g, f \rangle$):
\[
  \tr(\hat{O}_c^*)
  = \sum_\mu \langle s_\mu,\, \hat{O}_c^*\, s_\mu \rangle
  = \sum_\mu \langle \hat{O}_c\, s_\mu,\, s_\mu \rangle
  = \sum_\mu \langle s_\mu,\, \hat{O}_c\, s_\mu \rangle
  = \tr(\hat{O}_c). \qedhere
\]
\end{proof}

\begin{corollary}[Profile symmetry]
\label{cor:profile-symmetry}
By cyclic invariance of the trace, $\hat{O}_c^*$ can be rewritten
with $D_w$ at the end:
\begin{equation}\label{eq:O-adj-cyclic}
  \tr(\hat{O}_c^*)
  = \tr\!\bigl(
    G^{\bar\pi_{T-1}}(u_{T-1}) \cdots G^{\bar\pi_0}(u_0)\, D_w
  \bigr)
  = \GK_{\bar\pi^{\mathrm{rev}}}(\mathbf{u}^{\mathrm{rev}}, w).
\end{equation}
Combined with Proposition~\ref{prop:trace-invariance}, this gives:
\begin{equation}\label{eq:GK-symmetry}
  \boxed{%
  \GK_\pi(\mathbf{u}, w)
  = \GK_{\bar\pi^{\mathrm{rev}}}(\mathbf{u}^{\mathrm{rev}}, w).}
\end{equation}
The generating function for profile~$\pi$ equals that for the
reversed-complementary profile with reversed parameters.
\end{corollary}

\begin{remark}[The adjoint paradox]
\label{rem:paradox}
The adjoint is \textbf{invisible at the trace level}: $\hat{O}_c$
and $\hat{O}_c^*$ produce the same generating function.  The
adjoint becomes nontrivial only at the level of individual matrix
elements $\langle s_\lambda, \hat{O}_c\, s_\mu \rangle$.
\end{remark}


%% ============================================================
\section{The multivariate $E$--$H$ factorisation}
\label{sec:EH}
%% ============================================================

\subsection{Factorisation at the generating function level}

The growth diagram bijection
$\psi\colon (\IP, \ALCD(c)) \to \CPP(c)$ decomposes each cylindric
plane partition into a base partition $\gamma$ (contributing
$w^{|\gamma|}$ to every slice) and an ALCD $\mathfrak{d}$
(contributing $z_k^{|\mathrm{diag}(k)|_\mathfrak{d}}$ to the
$k$-th slice).  Setting $w = z_0 z_1 z_2$:
\begin{equation}\label{eq:EH-multi}
  \GK_c(\mathbf{z})
  = \underbrace{\frac{1}{(w;\,w)_\infty}}_{\displaystyle\sum_\gamma
    w^{|\gamma|}}
  \cdot\; E_c(\mathbf{z}),
\end{equation}
where
\begin{equation}\label{eq:Ec}
  E_c(\mathbf{z})
  := \sum_{\mathfrak{d} \in \ALCD(c)}
     \prod_{k=0}^{T-1} z_k^{|\mathrm{diag}(k)|_\mathfrak{d}}
\end{equation}
encodes the purely cylindric content.

The \emph{corrected multivariate generating function} is:
\begin{equation}\label{eq:Gc-multi}
  G_c(\mathbf{z})
  := (w;\,w)_\infty \cdot \GK_c(\mathbf{z})
  = E_c(\mathbf{z}).
\end{equation}

\subsection{The corrected operator}

\begin{definition}[Corrected operator]
\label{def:G-hat}
Define the \emph{corrected operator-valued generating function}:
\begin{equation}\label{eq:G-hat}
  \hat{G}_c(\mathbf{z})
  := (w;\,w)_\infty \cdot \hat{O}_c(\mathbf{z}),
\end{equation}
where $(w;\,w)_\infty = \prod_{k \geq 1}(1 - w^k)$ is a scalar
(depending on $w = z_0 z_1 z_2$) that multiplies the operator.
\end{definition}

Its trace recovers the corrected generating function:
\[
  \tr(\hat{G}_c(\mathbf{z}))
  = (w;\,w)_\infty \cdot \tr(\hat{O}_c(\mathbf{z}))
  = G_c(\mathbf{z}).
\]

\subsection{Multi-index coefficients}

Expand the corrected operator in a multi-index power series:
\begin{equation}\label{eq:G-hat-expand}
  \hat{G}_c(\mathbf{z})
  = \sum_{\mathbf{n} \in \N^T}
    \hat{G}_{c,\mathbf{n}}\;
    z_0^{n_0}\, z_1^{n_1} \cdots z_{T-1}^{n_{T-1}},
\end{equation}
where each $\hat{G}_{c,\mathbf{n}} \colon \La \to \La$ is an
operator.  Note that the Euler factor $(w;\,w)_\infty$ mixes
multi-index coefficients: $\hat{G}_{c,\mathbf{n}}$ is \emph{not}
simply $(w;\,w)_\infty$ times the $\mathbf{n}$-th coefficient
of $\hat{O}_c$, because $(w;\,w)_\infty$ depends on
$w = z_0 z_1 z_2$.

The trace of each coefficient gives the corrected multi-index
coefficient:
\begin{equation}\label{eq:gn-multi}
  g_{\mathbf{n},c}
  = \tr(\hat{G}_{c,\mathbf{n}})
  = [\mathbf{z}^\mathbf{n}]\, G_c(\mathbf{z}).
\end{equation}


%% ============================================================
\section{The $Q$-operator at Level~3}
\label{sec:Q-operator-L3}
%% ============================================================

\subsection{Definition}

\begin{definition}[$Q$-operator]
\label{def:Q-hat}
For a multi-index $\mathbf{n} = (n_0, n_1, \ldots, n_{T-1})
\in \N^T$, profile $c$ with $d \not\equiv 0 \pmod{3}$, define the
\emph{$Q$-operator}:
\begin{equation}\label{eq:Q-hat}
  \hat{Q}_{c,\mathbf{n}}
  := \prod_{k=0}^{T-1} \qpochn{q}{q}{n_k}
     \cdot \hat{G}_{c,\mathbf{n}}
  \;\colon\; \La \to \La,
\end{equation}
where $q = w = z_0 z_1 \cdots z_{T-1}$ and the
$\qpochn{q}{q}{n_k} = \prod_{j=1}^{n_k}(1-q^j)$ are scalar
factors that clear the denominators.
\end{definition}

\begin{remark}
In the rank-3 case, one often restricts to multi-indices of the
form $\mathbf{n} = (n_0, n_1, n_2) \in \N^3$ (one index per slice),
and the product becomes $\prod_{k=0}^{2} \qpochn{q}{q}{n_k}$.
\end{remark}

\subsection{Trace gives the multivariate $Q$-polynomial}

\begin{equation}\label{eq:Q-trace}
  Q_{\mathbf{n},c}
  = \tr(\hat{Q}_{c,\mathbf{n}})
  = \prod_{k} \qpochn{q}{q}{n_k} \cdot g_{\mathbf{n},c}.
\end{equation}

In terms of Schur basis matrix elements:
\begin{equation}\label{eq:Q-schur-sum}
  Q_{\mathbf{n},c}
  = \sum_\mu
    \langle s_\mu,\, \hat{Q}_{c,\mathbf{n}}\, s_\mu \rangle.
\end{equation}


%% ============================================================
\section{The adjoint $Q$-operator}
\label{sec:Q-adjoint}
%% ============================================================

\subsection{Definition}

\begin{definition}[Adjoint $Q$-operator]
\label{def:Q-hat-adj}
The adjoint is defined by:
\begin{equation}\label{eq:Q-adj-def}
  \langle \hat{Q}_{c,\mathbf{n}}^*\, f,\; g \rangle
  = \langle f,\; \hat{Q}_{c,\mathbf{n}}\, g \rangle
  \qquad \text{for all } f, g \in \La.
\end{equation}
\end{definition}

Since the $(q)$-Pochhammer factors are scalars:
\begin{equation}\label{eq:Q-adj-decomp}
  \hat{Q}_{c,\mathbf{n}}^*
  = \prod_{k} \qpochn{q}{q}{n_k} \cdot \hat{G}_{c,\mathbf{n}}^*.
\end{equation}

The nontrivial content is in $\hat{G}_{c,\mathbf{n}}^*$, which
inherits from $\hat{O}_c^*$ the profile transformation: the adjoint
corresponds to the \textbf{reversed-complementary profile}
$\bar\pi^{\mathrm{rev}}$ with reversed parameters.

\subsection{Matrix elements: the transpose relation}

In the Schur basis:
\begin{equation}\label{eq:Q-adj-matrix}
  \langle s_\lambda,\, \hat{Q}_{c,\mathbf{n}}^*\, s_\mu \rangle
  = \langle s_\mu,\, \hat{Q}_{c,\mathbf{n}}\, s_\lambda \rangle.
\end{equation}
The matrix of $\hat{Q}_{c,\mathbf{n}}^*$ is the \emph{transpose}
of the matrix of $\hat{Q}_{c,\mathbf{n}}$.

\subsection{Trace invariance}

\begin{proposition}
\label{prop:Q-trace-invariance}
$\tr(\hat{Q}_{c,\mathbf{n}}^*) = \tr(\hat{Q}_{c,\mathbf{n}})
= Q_{\mathbf{n},c}$.
\end{proposition}

\begin{proof}
\[
  \tr(\hat{Q}_{c,\mathbf{n}}^*)
  = \prod_k \qpochn{q}{q}{n_k} \cdot \tr(\hat{G}_{c,\mathbf{n}}^*)
  = \prod_k \qpochn{q}{q}{n_k} \cdot \tr(\hat{G}_{c,\mathbf{n}})
  = Q_{\mathbf{n},c},
\]
using Proposition~\ref{prop:trace-invariance} and the fact that
scalar factors commute with the adjoint.
\end{proof}

\begin{remark}[The adjoint is invisible at the polynomial level]
The adjoint $\hat{Q}^*_{c,\mathbf{n}}$ produces the same
cylindric Kostka polynomial as $\hat{Q}_{c,\mathbf{n}}$.
Its content is visible only at the level of individual matrix
elements (the kernel).
\end{remark}


%% ============================================================
\section{The kernel and the positivity conjecture}
\label{sec:kernel}
%% ============================================================

\subsection{The $Q$-kernel}

\begin{definition}
The \emph{$Q$-kernel} is:
\begin{equation}\label{eq:kernel}
  K_{c,\mathbf{n}}(\lambda, \mu)
  := \langle s_\lambda,\, \hat{Q}_{c,\mathbf{n}}\, s_\mu \rangle
  \;\colon\; \IP \times \IP \to \Q[q].
\end{equation}
\end{definition}

The adjoint kernel is the transpose:
$K_{c,\mathbf{n}}^*(\lambda, \mu) = K_{c,\mathbf{n}}(\mu, \lambda)$.
The diagonal is unchanged: $K^*(\mu,\mu) = K(\mu,\mu)$.

\subsection{The positivity conjecture}

\begin{conjecture}[Multivariate positivity, operator form]
\label{conj:positivity}
For all $\mathbf{n} \in \N^T$ and all profiles $c$ with
$d \not\equiv 0 \pmod{3}$:
\begin{equation}\label{eq:positivity}
  \tr(\hat{Q}_{c,\mathbf{n}})
  = \sum_\mu K_{c,\mathbf{n}}(\mu, \mu)
  \in \Z_{\geq 0}[q].
\end{equation}
\end{conjecture}

This is a statement about the \textbf{sum of diagonal matrix
elements}.  The conjecture does not require:
\begin{enumerate}
  \item individual diagonal elements
    $K_{c,\mathbf{n}}(\mu, \mu)$ to be nonneg,
  \item off-diagonal elements
    $K_{c,\mathbf{n}}(\lambda, \mu)$ to have any sign, or
  \item self-adjointness ($\hat{Q} = \hat{Q}^*$).
\end{enumerate}

\subsection{The symmetrised operator}

The \emph{symmetrised} $Q$-operator
\begin{equation}\label{eq:symmetrised}
  \hat{Q}_{c,\mathbf{n}}^{\mathrm{sym}}
  := \tfrac{1}{2}(\hat{Q}_{c,\mathbf{n}}
     + \hat{Q}_{c,\mathbf{n}}^*)
\end{equation}
is self-adjoint by construction, with the same trace:
$\tr(\hat{Q}^{\mathrm{sym}}) = Q_{\mathbf{n},c}$.

Its kernel is symmetric:
$K^{\mathrm{sym}}(\lambda, \mu)
= \tfrac{1}{2}(K(\lambda,\mu) + K(\mu,\lambda))$.
When restricted to a finite-dimensional subspace of $\La$, the
eigenvalues of $\hat{Q}^{\mathrm{sym}}$ are real, and their sum
equals $Q_{\mathbf{n},c}$.

\begin{remark}[Positive semi-definiteness]
If $\hat{Q}^{\mathrm{sym}}$ were positive semi-definite (all
eigenvalues $\geq 0$), then $Q_{\mathbf{n},c} \geq 0$.  This is
strictly stronger than the conjecture (it implies each diagonal
element is nonneg, not just their sum).
\end{remark}


%% ============================================================
\section{The commutation relations at Level~3}
\label{sec:commutation-L3}
%% ============================================================

\subsection{Vertex operator commutation}

The fundamental commutation relation (Schur case) is:
\begin{equation}\label{eq:comm}
  \Omega^*[Xu]\;\Omega[Xv]
  = \frac{1}{1 - uv}\;\Omega[Xv]\;\Omega^*[Xu].
\end{equation}
Each commutation produces a scalar factor $1/(1-uv)
= \sum_{m \geq 0}(uv)^m$, generating the face labels of the
growth diagram.

\subsection{Degree operator commutation}

\begin{align}
  D_w\, \Omega[Xu] &= \Omega[Xuw]\, D_w,
    \label{eq:Dw-comm-up} \\
  D_w\, \Omega^*[Xu] &= \Omega^*[Xu/w]\, D_w.
    \label{eq:Dw-comm-down}
\end{align}

These are used in the proof of the Borodin product formula:
commuting all $\Omega^*$ past all $\Omega$ produces scalar
factors, and iterating the $D_w$ shift produces the infinite
product.

\subsection{Self-commutation}

Since $(\Omega[Xv])^* = \Omega^*[Xv]$, the fundamental relation
can be read as:
\begin{equation}\label{eq:self-comm}
  (\Omega[Xv])^*\;\Omega[Xv]
  = \frac{1}{1 - v^2}\;\Omega[Xv]\;(\Omega[Xv])^*.
\end{equation}
The scalar factor $1/(1-v^2)$ measures the non-commutativity of
$\Omega[Xv]$ with its own adjoint.


%% ============================================================
\section{From multivariate to bivariate}
\label{sec:multi-to-bi}
%% ============================================================

The bivariate generating function $\GK_c(z,q) = \sum_\mathfrak{c}
z^{\max(\mathfrak{c})} q^{|\mathfrak{c}|}$ is obtained from the
multivariate one by the non-linear \emph{max extraction}:
\[
  z_0^{|\mu^0|}\, z_1^{|\mu^1|}\, z_2^{|\mu^2|}
  \;\longmapsto\;
  z^{\max_k \mu^k_1}\, q^{|\mathfrak{c}|}.
\]

This does not correspond to a specialisation of the $z_k$, so it
does not lift to the operator level.  The bivariate $Q_{n,c}(q)$
cannot be obtained from $\hat{Q}_{c,\mathbf{n}}$ by a simple
specialisation.  The chain of implications is:
\[
  \text{Conjecture~\ref{conj:positivity}
  (multivariate)}
  \;\;\underset{\text{max projection}}{\Longrightarrow}\;\;
  \text{Bivariate positivity (Warnaar, Conj.~2.7)}.
\]


%% ============================================================
\section{Summary of Part~I}
\label{sec:summary-I}
%% ============================================================

\begin{center}
\begin{tabular}{l|l|l}
\textbf{Object} & \textbf{Definition} & \textbf{Role} \\ \hline
$\hat{O}_c$ &
  $\prod_k G^{\pi_k}(u_k)\, D_w$ &
  Level~3 operator on $\La$ \\[4pt]
$\hat{O}_c^*$ &
  $D_w \prod_k G^{\bar\pi_{T-1-k}}(u_{T-1-k})$ &
  Reversed-complementary profile \\[4pt]
$\hat{G}_c(\mathbf{z})$ &
  $(w;w)_\infty \cdot \hat{O}_c(\mathbf{z})$ &
  Corrected operator ($E$--$H$) \\[4pt]
$\hat{G}_{c,\mathbf{n}}$ &
  $[\mathbf{z}^\mathbf{n}]\,\hat{G}_c(\mathbf{z})$ &
  Multi-index coefficient operator \\[4pt]
$\hat{Q}_{c,\mathbf{n}}$ &
  $\prod_k (q;q)_{n_k} \cdot \hat{G}_{c,\mathbf{n}}$ &
  $Q$-operator (denominators cleared) \\[4pt]
$\hat{Q}_{c,\mathbf{n}}^*$ &
  $\prod_k (q;q)_{n_k} \cdot \hat{G}_{c,\mathbf{n}}^*$ &
  Adjoint $Q$-operator \\[4pt]
$K_{c,\mathbf{n}}(\lambda,\mu)$ &
  $\langle s_\lambda, \hat{Q}_{c,\mathbf{n}}\, s_\mu \rangle$ &
  $Q$-kernel \\[4pt]
$Q_{\mathbf{n},c}$ &
  $\sum_\mu K_{c,\mathbf{n}}(\mu,\mu) = \tr(\hat{Q}_{c,\mathbf{n}})$ &
  Multivariate cylindric Kostka poly.
\end{tabular}
\end{center}

\medskip
Key structural facts:
\begin{enumerate}
  \item $\hat{O}_c^*$ corresponds to the reversed-complementary
    profile $\bar\pi^{\mathrm{rev}}$.
  \item $\tr(\hat{Q}^*) = \tr(\hat{Q}) = Q_{\mathbf{n},c}$: the
    adjoint is invisible at the polynomial level.
  \item The positivity conjecture is a statement about the trace
    (sum of diagonal matrix elements).
  \item $\hat{Q}^{\mathrm{sym}}
    = \frac{1}{2}(\hat{Q} + \hat{Q}^*)$ is self-adjoint with the
    same trace; positive semi-definiteness would imply the
    conjecture but is strictly stronger.
  \item The max extraction producing the bivariate $Q_{n,c}(q)$
    does not lift to the operator level.
\end{enumerate}


%% ============================================================
%% PART II: THE BIVARIATE Q-OPERATOR AND CASCADE COMMUTATORS
%% ============================================================

\part{The Bivariate $Q$-Operator and Cascade Commutators}


%% ============================================================
\section{The operators $P_m$ and $T_n$}
\label{sec:PT}
%% ============================================================

We use the following mode operators on $\La$:
\begin{itemize}
  \item $P_m = h_m \cdot{}$: multiplication by the complete
    homogeneous symmetric function $h_m$.  In generating function
    form, $P(z) = \sum_{m \geq 0} P_m z^m = \Omega[Xz]$.
  \item $T_n = P_n^* = \partial_{h_n}$: the adjoint of $P_n$ with
    respect to the Hall inner product.  In generating function form,
    $T(z) = \sum_{n \geq 0} T_n z^n$.
\end{itemize}

These satisfy the fundamental commutation relation
(cf.~\cite{langer-masters}, \cite{COMMUTATION}):
\begin{equation}\label{eq:fund-comm}
  T(v)\, P(w) = \frac{1}{1 - vw}\; P(w)\, T(v).
\end{equation}


%% ============================================================
\section{The bivariate $\mathbf{Q}(z)$ and $\mathbf{Q}^*(w)$}
\label{sec:Q-bivariate}
%% ============================================================

\begin{definition}\label{def:Q-bi}
The \emph{$Q$-operator} (bivariate version) for profile $c$ is the
formal power series in $z$ with operator coefficients:
\begin{equation}\label{eq:Q-bi}
  \mathbf{Q}(z)
  = \sum_{m \geq 0} Q_{m,c}(q)\, z^m\, P_m
  = \sum_{m \geq 0} Q_{m,c}(q)\, z^m\, h_m \cdot{},
\end{equation}
where $Q_{0,c}(q) = 1$.

The \emph{adjoint $Q$-operator} is:
\begin{equation}\label{eq:Qstar-bi}
  \mathbf{Q}^*(w)
  = \sum_{n \geq 0} Q_{n,c}(q)\, w^n\, T_n
  = \sum_{n \geq 0} Q_{n,c}(q)\, w^n\, \partial_{h_n}.
\end{equation}
\end{definition}

\begin{remark}
This is a fundamentally different construction from
Part~I.  In Part~I, the $Q$-operator $\hat{Q}_{c,\mathbf{n}}$
is a single operator whose \emph{trace} gives the polynomial
$Q_{\mathbf{n},c}$.  Here, the polynomials $Q_m(q)$ appear as
\emph{coefficients} multiplying the Pieri operators $P_m$ and their
adjoints $T_n$.
\end{remark}


%% ============================================================
\section{Commutation relations in the $h$-basis}
\label{sec:h-comm}
%% ============================================================

\subsection{Mode-level commutation}

Extracting the coefficient of $v^n w^m$ from~\eqref{eq:fund-comm}:

\begin{proposition}\label{prop:TnPm}
\begin{equation}\label{eq:TnPm}
  T_n\, P_m = \sum_{k=0}^{\min(n,m)} P_{m-k}\, T_{n-k}.
\end{equation}
\end{proposition}

\begin{proof}
From $T(v)\, P(w) = \frac{1}{1-vw}\, P(w)\, T(v)$, the right-hand
side is
\[
  \Bigl(\sum_{k \geq 0} (vw)^k\Bigr)
  \Bigl(\sum_a P_a\, w^a\Bigr)
  \Bigl(\sum_b T_b\, v^b\Bigr)
  = \sum_{k,a,b} P_a\, T_b\, v^{k+b}\, w^{k+a}.
\]
Setting $n = k + b$ and $m = k + a$ gives $k$ ranging from $0$
to $\min(n,m)$, with $a = m - k$ and $b = n - k$.
\end{proof}

\begin{corollary}\label{cor:commutator}
The commutator is:
\begin{equation}\label{eq:comm-TnPm}
  [T_n, P_m] = \sum_{k=1}^{\min(n,m)} P_{m-k}\, T_{n-k}.
\end{equation}
In particular:
\begin{align}
  [T_1, P_1] &= \mathrm{id},  \label{eq:T1P1} \\
  [T_n, P_n] &= P_{n-1} T_{n-1} + P_{n-2} T_{n-2}
    + \cdots + \mathrm{id}.  \label{eq:TnPn}
\end{align}
\end{corollary}

\begin{remark}[Comparison with power sum basis]
In the power sum basis, $[\frac{1}{n}\partial_{p_n},\, p_m \cdot]
= \delta_{n,m}$ (the Heisenberg algebra).  The $h$-basis introduces
\emph{cascade terms} $P_{m-k} T_{n-k}$ for intermediate values
of~$k$.  These cross-level terms are the key structural feature.
\end{remark}


\subsection{Lowering applied to $h_m$}

\begin{lemma}\label{lem:Tn-hm}
\begin{equation}\label{eq:Tn-hm}
  T_n(h_m) =
  \begin{cases}
    h_{m-n} & \text{if } m \geq n, \\
    0 & \text{if } m < n.
  \end{cases}
\end{equation}
\end{lemma}

\begin{proof}
$T_n(h_m) = T_n P_m(1) = \sum_{k=0}^{\min(n,m)} P_{m-k} T_{n-k}(1)$.
Since $T_j(1) = 0$ for $j \geq 1$, only $k = n$ survives
(when $n \leq m$), giving $P_{m-n}(1) = h_{m-n}$.
\end{proof}


%% ============================================================
\section{Operator products applied to the vacuum}
\label{sec:products}
%% ============================================================

\subsection{The product $\mathbf{Q}^*(w)\,\mathbf{Q}(z)$}

\begin{proposition}\label{prop:QstarQ}
\begin{equation}\label{eq:QstarQ-vacuum}
  \mathbf{Q}^*(w)\,\mathbf{Q}(z)\,[1]
  = \sum_{a \geq 0} C_a(zw)\, z^a\, h_a,
\end{equation}
where $C_a(t) := \sum_{n \geq 0} Q_{a+n}(q)\, Q_n(q)\, t^n$ is
the cross-correlation function.
\end{proposition}

\begin{proof}
By~\eqref{eq:Q-bi},
$\mathbf{Q}(z)[1] = \sum_m Q_m z^m h_m$.  Applying $\mathbf{Q}^*(w)$:
\begin{align*}
  \mathbf{Q}^*(w)\bigl[\sum_m Q_m z^m h_m\bigr]
  &= \sum_{m,n} Q_m Q_n\, z^m w^n\, T_n(h_m) \\
  &= \sum_{\substack{m,n \\ n \leq m}} Q_m Q_n\, z^m w^n\, h_{m-n}
    \qquad\text{(by Lemma~\ref{lem:Tn-hm})} \\
  &= \sum_{a \geq 0} h_a\, z^a \sum_{n \geq 0}
    Q_{a+n}\, Q_n\, (zw)^n,
\end{align*}
substituting $a = m - n$.
\end{proof}


\subsection{The product $\mathbf{Q}(z)\,\mathbf{Q}^*(w)$}

Since $T_n(1) = 0$ for $n \geq 1$:
\begin{equation}\label{eq:QQstar-vacuum}
  \mathbf{Q}(z)\,\mathbf{Q}^*(w)\,[1]
  = \mathbf{Q}(z)\,[1]
  = \sum_{m \geq 0} Q_m\, z^m\, h_m.
\end{equation}


\subsection{The commutator on the vacuum}

\begin{theorem}\label{thm:commutator-vacuum}
\begin{equation}\label{eq:comm-vacuum}
  [\mathbf{Q}(z),\, \mathbf{Q}^*(w)]\,[1]
  = -\sum_{a \geq 0} z^a\, h_a
    \sum_{n \geq 1} Q_{a+n}\, Q_n\, (zw)^n.
\end{equation}
\end{theorem}

\begin{proof}
Subtracting~\eqref{eq:QstarQ-vacuum}
from~\eqref{eq:QQstar-vacuum}:
\begin{align*}
  [\mathbf{Q}(z), \mathbf{Q}^*(w)][1]
  &= \sum_a Q_a z^a h_a
    - \sum_a z^a h_a \sum_{n \geq 0} Q_{a+n} Q_n (zw)^n \\
  &= \sum_a z^a h_a \bigl(Q_a - Q_a Q_0
    - \textstyle\sum_{n \geq 1} Q_{a+n} Q_n (zw)^n\bigr).
\end{align*}
Since $Q_0 = 1$, the first two terms cancel.
\end{proof}


%% ============================================================
\section{The vacuum expectation value}
\label{sec:vev}
%% ============================================================

Taking the Hall inner product of~\eqref{eq:QstarQ-vacuum} with $1$
(extracting the $h_0$ coefficient):

\begin{corollary}\label{cor:vev}
\begin{equation}\label{eq:vev}
  \langle 1 \mid \mathbf{Q}^*(w)\,\mathbf{Q}(z) \mid 1 \rangle
  = C_0(zw) = \sum_{n \geq 0} Q_n(q)^2\, (zw)^n.
\end{equation}
\end{corollary}

\begin{remark}
If each $Q_n(q) \in \N[q]$ (the positivity conjecture), then
$Q_n(q)^2 \in \N[q]$ and the vacuum expectation value is manifestly
$q$-positive.  The converse does not hold: $Q_n^2 \in \N[q]$ does
not imply $Q_n \in \N[q]$.
\end{remark}


%% ============================================================
\section{The cross-correlation functions}
\label{sec:cross-corr}
%% ============================================================

\begin{definition}\label{def:Ca}
The \emph{cross-correlation functions} are:
\begin{equation}\label{eq:Ca}
  C_a(t) = \sum_{n \geq 0} Q_{a+n}(q)\, Q_n(q)\, t^n
  \qquad (a \geq 0).
\end{equation}
\end{definition}

\subsection{$d = 2$ (Rogers--Ramanujan)}

For $c = (1,1,0)$, $Q_n = q^{n^2}$, so:
\begin{equation}\label{eq:Ca-RR}
  C_a(t) = q^{a^2} \sum_{n \geq 0} q^{2n^2 + 2an}\, t^n.
\end{equation}
The vacuum expectation is the theta-like sum
$C_0(t) = \sum_{n \geq 0} q^{2n^2}\, t^n$.


\subsection{$d = 4$}

For $c = (2,2,0)$, the first few polynomials are:
\begin{align*}
  Q_0 &= 1, \\
  Q_1 &= q^2 + q^3 + q^4 + q^5.
\end{align*}
The vacuum expectation begins:
$C_0(t) = 1 + (q^4 + 2q^5 + 3q^6 + 4q^7 + 3q^8
+ 2q^9 + q^{10})\, t + \cdots$


%% ============================================================
\section{The full commutator on $\La$}
\label{sec:full-comm}
%% ============================================================

The full commutator of the $Q$-operators acting on a general
symmetric function $f \in \La$ takes the form:

\begin{theorem}\label{thm:full-comm}
\begin{equation}\label{eq:full-comm}
  [\mathbf{Q}(z), \mathbf{Q}^*(w)]
  = -\sum_{a,b \geq 0} P_a\, T_b
    \sum_{k \geq 1} Q_{a+k}\, Q_{b+k}\, z^{a+k}\, w^{b+k}.
\end{equation}
\end{theorem}

\begin{proof}
Using the commutation $T_n P_m = \sum_{j=0}^{\min(n,m)} P_{m-j}
T_{n-j}$ from Proposition~\ref{prop:TnPm}:
\begin{align*}
  \mathbf{Q}^*(w)\, \mathbf{Q}(z)
  &= \sum_{m,n} Q_m Q_n\, z^m w^n\, T_n P_m \\
  &= \sum_{m,n} Q_m Q_n\, z^m w^n
    \sum_{j=0}^{\min(n,m)} P_{m-j}\, T_{n-j}.
\end{align*}
Setting $a = m - j$, $b = n - j$, $k = j$:
\[
  \mathbf{Q}^*(w)\, \mathbf{Q}(z)
  = \sum_{a,b \geq 0} P_a\, T_b\; z^a w^b
    \sum_{k \geq 0} Q_{a+k}\, Q_{b+k}\, (zw)^k.
\]
The $\mathbf{Q}(z)\, \mathbf{Q}^*(w)$ product has only the
$k = 0$ term (since $P_a$ and $T_b$ act independently):
\[
  \mathbf{Q}(z)\, \mathbf{Q}^*(w)
  = \sum_{a,b \geq 0} Q_a\, Q_b\; z^a w^b\; P_a\, T_b.
\]
Subtracting gives the commutator with the $k = 0$ terms cancelling.
\end{proof}

\begin{remark}[Cascade structure]
The coefficient of $P_a T_b$ in the commutator is
\[
  -z^a w^b \sum_{k \geq 1} Q_{a+k}(q)\, Q_{b+k}(q)\, (zw)^k,
\]
involving \emph{shifted} cross-correlations.  The presence of terms
with $a \neq b$ is the key difference from the power sum basis
formulation, where the commutator is diagonal:
$[\sum Q_m z^m p_m,\, \sum Q_n w^n \frac{1}{n}\partial_{p_n}]
= \sum_{n \geq 1} Q_n^2 (zw)^n$.
\end{remark}

\begin{remark}[Cross-level coupling and the involution]
The $P_a T_b$ terms with $a \neq b$ couple different depth levels of
the cylindric partition.  In the involution framework
(see~\cite{FIRST_ATTEMPT}), the involution must pair elements at
level~$m$ with elements at level~$m'$.  The cascade terms provide
the algebraic mechanism for such cross-level interactions.
\end{remark}


%% ============================================================
%% PART III: THE GENERALIZED LOCAL RULE
%% ============================================================

\part{The Generalized Local Rule}


%% ============================================================
\section{The standard algebra--bijection correspondence}
\label{sec:std-correspondence}
%% ============================================================

We begin by recalling the precise dictionary between the algebraic
framework (commutation relations) and the bijective framework
(growth diagram local rules).  This dictionary is the foundation on
which the $Q$-enriched generalisation is built.

\subsection{The Burge local rule as a commutation step}

The fundamental commutation relation~\eqref{eq:fund-comm},
\[
  T(v)\, P(w) = \frac{1}{1 - vw}\; P(w)\, T(v),
\]
is \emph{equivalent} to the Burge local rule at the mode level.
Extracting the coefficient of $v^n w^m$ from both sides gives the
cascade relation (Proposition~\ref{prop:TnPm}):
\begin{equation}\label{eq:cascade-recap}
  T_n\, P_m = \sum_{k=0}^{\min(n,m)} P_{m-k}\, T_{n-k}.
\end{equation}
Each term in the sum corresponds to one possible output of the
local rule: given an input diamond $(\nu;\, \alpha, \beta)$ where
$|\alpha| - |\nu| = m$ and $|\beta| - |\nu| = n$, the cascade
parameter $k$ determines the output
$\lambda = \mathfrak{U}_{\alpha,\beta}(k,\nu)$ with
$|\lambda| = |\alpha| + |\beta| - |\nu| + k$.

\begin{remark}[The dictionary]
\label{rem:dictionary}
The correspondence between the algebraic proof and the bijective
proof is summarised as follows
(cf.~\cite{COMMUTATION}, Remark~8):
\begin{center}
\begin{tabular}{l|l}
\textbf{Algebraic operation} & \textbf{Bijective operation} \\
\hline
Commuting $T(v)$ past $P(w)$ &
  Applying the local rule $\mathfrak{U}_{\alpha,\beta}$ \\[3pt]
Scalar factor $\frac{1}{1-vw} = \sum_k (vw)^k$ &
  Face label $k \in \Z_{\geq 0}$ \\[3pt]
Cascade parameter $k$ &
  Overlap $|A \cap B|$ in conjugate coordinates \\[3pt]
Cyclic invariance of trace &
  Rotation of the cylinder \\[3pt]
Commuting all $T$'s past all $P$'s &
  Full growth diagram computation \\[3pt]
$D_w$ shift: $D_w P(u) = P(uw) D_w$ &
  Depth increment (adding base partition)
\end{tabular}
\end{center}
\end{remark}

\subsection{Face labels and the generating function}

In the standard growth diagram for profile~$c$, the face labels
are non-negative integers $k \in \Z_{\geq 0}$, each contributing
weight~$1$ regardless of the value of~$k$.  The face-label
generating function for a single diamond is therefore:
\begin{equation}\label{eq:face-gf-std}
  \sum_{k \geq 0} t^k = \frac{1}{1-t},
\end{equation}
which is precisely the scalar factor arising from commuting $T(v)$
past $P(w)$ with $t = vw$.  This is \emph{universal}: it depends
neither on the profile~$c$ nor on the polynomials~$Q_n(q)$.


%% ============================================================
\section{The $Q$-kernel}
\label{sec:Q-kernel}
%% ============================================================

\subsection{Derivation from the cascade commutator}

The cascade commutator (Theorem~\ref{thm:full-comm}) showed that
the product $\mathbf{Q}^*(w)\, \mathbf{Q}(z)$ takes the form:
\begin{equation}\label{eq:QstarQ-full}
  \mathbf{Q}^*(w)\, \mathbf{Q}(z)
  = \sum_{a,b \geq 0} P_a\, T_b\; z^a w^b
    \sum_{k \geq 0} Q_{a+k}\, Q_{b+k}\, (zw)^k.
\end{equation}
This motivates extracting the inner sum as an independent object.

\begin{definition}[$Q$-kernel]
\label{def:Phi}
The \emph{$Q$-kernel} is the bivariate function:
\begin{equation}\label{eq:Phi}
  \Phi_{a,b}(t)
  := \sum_{k \geq 0} Q_{a+k}(q)\, Q_{b+k}(q)\, t^k
  \qquad (a, b \geq 0),
\end{equation}
where $Q_n = Q_{n,c}(q)$ are the cylindric Kostka polynomials for
profile~$c$.
\end{definition}

With this notation, the cascade product becomes:
\begin{equation}\label{eq:QstarQ-Phi}
  \mathbf{Q}^*(w)\, \mathbf{Q}(z)
  = \sum_{a,b \geq 0} P_a\, T_b\; z^a w^b\; \Phi_{a,b}(zw),
\end{equation}
and the commutator (Theorem~\ref{thm:full-comm}) isolates the
$k \geq 1$ terms:
\begin{equation}\label{eq:comm-Phi}
  [\mathbf{Q}(z),\, \mathbf{Q}^*(w)]
  = -\sum_{a,b \geq 0} P_a\, T_b\; z^a w^b\;
    \bigl(\Phi_{a,b}(zw) - Q_a Q_b\bigr).
\end{equation}

\subsection{Relation to the cross-correlations}

The $Q$-kernel unifies the cross-correlation functions from
Section~\ref{sec:cross-corr}:

\begin{proposition}\label{prop:Phi-Ca}
\begin{equation}\label{eq:Phi-Ca}
  \Phi_{a,0}(t) = C_a(t) = \sum_{n \geq 0} Q_{a+n}\, Q_n\, t^n
  \qquad (a \geq 0).
\end{equation}
In particular, $\Phi_{0,0}(t) = C_0(t) = \sum_{n \geq 0}
Q_n^2\, t^n$ is the vacuum expectation value.
\end{proposition}

\begin{proof}
Set $b = 0$ in~\eqref{eq:Phi} and use $Q_0 = 1$.
\end{proof}

\subsection{Properties of the $Q$-kernel}

\begin{proposition}\label{prop:Phi-properties}
The $Q$-kernel satisfies:
\begin{enumerate}
  \item \emph{Initial value.}\;
    $\Phi_{a,b}(0) = Q_a\, Q_b$.
  \item \emph{Symmetry.}\;
    $\Phi_{a,b}(t) = \Phi_{b,a}(t)$.
  \item \emph{Vacuum.}\;
    $\Phi_{0,0}(t) = C_0(t)$.
  \item \emph{Cross-correlation.}\;
    $\Phi_{a,0}(t) = C_a(t)$.
  \item \emph{Shifted factorisation.}\;
    $\Phi_{a,b}(t) = Q_a Q_b + t\, \Phi_{a+1,b+1}(t)$.
\end{enumerate}
\end{proposition}

\begin{proof}
Properties (1)--(4) follow directly from the definition.
For~(5), separate the $k = 0$ term:
\[
  \Phi_{a,b}(t)
  = Q_a Q_b + \sum_{k \geq 1} Q_{a+k} Q_{b+k}\, t^k
  = Q_a Q_b + t \sum_{k \geq 0} Q_{a+1+k} Q_{b+1+k}\, t^k
  = Q_a Q_b + t\, \Phi_{a+1,b+1}(t). \qedhere
\]
\end{proof}

\begin{remark}[The $Q$-kernel as a deformation of $1/(1-t)$]
\label{rem:deformation}
When $Q_n = 1$ for all~$n$ (the $d = 1$ case),
$\Phi_{a,b}(t) = \sum_{k \geq 0} t^k = 1/(1-t)$, independent
of~$a$ and~$b$.  For general profiles, $\Phi_{a,b}(t)$ is a
\emph{profile-dependent} deformation of this universal kernel.
The shifted factorisation in property~(5) can be iterated:
\[
  \Phi_{a,b}(t)
  = \sum_{j=0}^{N-1} Q_{a+j}\, Q_{b+j}\, t^j
    + t^N \Phi_{a+N, b+N}(t),
\]
so the $Q$-kernel is determined by the sequence
$(Q_{a+k} Q_{b+k})_{k \geq 0}$ in the same way that $1/(1-t)$ is
determined by the constant sequence $(1, 1, 1, \ldots)$.
\end{remark}


%% ============================================================
\section{The generalised local rule}
\label{sec:gen-local-rule}
%% ============================================================

\subsection{Definition}

The standard Burge local rule assigns to each input diamond
$(\nu;\, \alpha, \beta)$ and cascade parameter $k \geq 0$ an
output partition $\lambda = \mathfrak{U}_{\alpha,\beta}(k,\nu)$
with weight $|\lambda| = |\alpha| + |\beta| - |\nu| + k$.
In the growth diagram, the cascade parameter becomes the face
label, and each label carries uniform weight~$1$.

The $Q$-enriched local rule uses the \emph{same} partition
geometry but assigns \emph{non-uniform} weights to the cascade
parameter.

\begin{definition}[$Q$-enriched local rule]
\label{def:Q-local-rule}
Given an input diamond $(\nu;\, \alpha, \beta)$ with
$|\alpha| - |\nu| = a$ and $|\beta| - |\nu| = b$, the
\emph{$Q$-enriched local rule} assigns to the cascade
parameter~$k$ the weight:
\begin{equation}\label{eq:Q-face-weight}
  \mathrm{wt}(k;\, a, b) = Q_{a+k}(q) \cdot Q_{b+k}(q).
\end{equation}
The output partition is the same as in the standard rule:
$\lambda = \mathfrak{U}_{\alpha,\beta}(k,\nu)$, with
$|\lambda| = |\alpha| + |\beta| - |\nu| + k$.
\end{definition}

The face-label generating function for the $Q$-enriched rule is
therefore:
\begin{equation}\label{eq:face-gf-Q}
  \sum_{k \geq 0} \mathrm{wt}(k;\, a, b)\, t^k
  = \sum_{k \geq 0} Q_{a+k}\, Q_{b+k}\, t^k
  = \Phi_{a,b}(t).
\end{equation}
This replaces the universal $1/(1-t)$ of the standard rule with
the profile-dependent $Q$-kernel.

\begin{remark}[Weight depends on context]
\label{rem:context-weight}
In the standard rule, the face-label generating function $1/(1-t)$
is the same for every diamond, regardless of the strip sizes $a$
and~$b$.  In the $Q$-enriched rule, the generating function
$\Phi_{a,b}(t)$ depends on the strip sizes.  This means the
$Q$-enriched rule is \emph{context-sensitive}: the weight assigned
to a face label depends on the horizontal and vertical strips
entering the diamond.
\end{remark}

\subsection{Weight preservation}

The $Q$-enriched rule preserves the weight
condition~\eqref{eq:cascade-recap} of the standard rule.

\begin{proposition}\label{prop:weight-preservation}
Let $(\nu;\, \alpha, \beta)$ be an input diamond with
$|\alpha| - |\nu| = a$ and $|\beta| - |\nu| = b$.  For each
cascade parameter $k \geq 0$, the output
$\lambda = \mathfrak{U}_{\alpha,\beta}(k,\nu)$ satisfies:
\begin{equation}\label{eq:weight-balance}
  |\lambda| = |\alpha| + |\beta| - |\nu| + k = a + b + |\nu| + k.
\end{equation}
This is a property of the Burge local rule itself and is
independent of the weight assignment.
\end{proposition}

\subsection{Invertibility}

\begin{proposition}\label{prop:Q-invertible}
The $Q$-enriched local rule is invertible: given
$(\lambda;\, \alpha, \beta)$, the cascade parameter $k$ is
uniquely determined by
$k = |\lambda| - |\alpha| - |\beta| + |\nu|$, and the inner
partition $\nu$ is recovered by the inverse Burge rule
$\nu = \mathfrak{U}_{\alpha,\beta}^{-1}(k, \lambda)$.
\end{proposition}

\begin{proof}
Invertibility is a property of the underlying Burge local rule
(see~\cite{UNDERLYING_SET}, Proposition~2.20): the map
$(\nu, k) \mapsto \lambda$ is a bijection from
$\{(\nu, k) : \nu \in D(\alpha) \cap D(\beta),\; k \geq 0\}$
to $U(\alpha) \cap U(\beta)$.  The $Q$-enriched rule uses the
same bijection; only the weight assignment changes.
\end{proof}


%% ============================================================
\section{The $Q$-enriched growth diagram}
\label{sec:Q-growth-diagram}
%% ============================================================

\subsection{Standard growth diagrams and ALCDs}

In the standard growth diagram for profile~$c$, one fills a
cylindric grid with partitions using the Burge local rule at each
face.  The face labels are non-negative integers, and the
collection of all face labels forms an Arbitrarily Labelled
Cylindric Diagram ($\ALCD$).  The key theorem
(cf.~\cite{UNDERLYING_SET}, Theorem~2.21) is that this gives a
weight-preserving bijection:
\begin{equation}\label{eq:bijection-std}
  \psi_c \colon (\IP,\, \ALCD(c))
  \;\xrightarrow{\;\sim\;}\;
  \CPP(c),
\end{equation}
where $\IP$ denotes the set of all integer partitions (the base
partition $\gamma$).  The multivariate generating function is:
\begin{equation}\label{eq:GK-std}
  \GK_c(\mathbf{z})
  = \frac{1}{(w;\,w)_\infty} \cdot
    \sum_{\mathfrak{d} \in \ALCD(c)}
    \prod_{k=0}^{T-1} z_k^{|\mathrm{diag}(k)|_\mathfrak{d}}.
\end{equation}

\subsection{The $Q$-enriched growth diagram}

\begin{definition}[$Q$-ALCD]
\label{def:Q-ALCD}
A \emph{$Q$-weighted Arbitrarily Labelled Cylindric Diagram}
($Q$-$\ALCD$) is an $\ALCD$ $\mathfrak{d}$ together with the
weight assignment:
\begin{equation}\label{eq:Q-ALCD-weight}
  \mathrm{wt}_Q(\mathfrak{d})
  := \prod_{\text{faces } f}
     Q_{a_f + \mathfrak{d}(f)}(q) \cdot Q_{b_f + \mathfrak{d}(f)}(q),
\end{equation}
where for each face $f$, the values $a_f$ and $b_f$ are the
horizontal and vertical strip sizes entering that face.
\end{definition}

The $Q$-enriched growth diagram uses the same partition geometry as
the standard one---the same Burge local rule, the same set of
valid $\ALCD$s---but assigns the $Q$-dependent weight to each
face rather than the uniform weight~$1$.

\begin{remark}[Two levels of growth diagram]
\label{rem:two-levels}
This gives a natural hierarchy:
\begin{center}
\begin{tabular}{lll}
\textbf{Level} & \textbf{Face weight} &
  \textbf{Generating function} \\[2pt]
\hline
Standard & $1$ (uniform) &
  Borodin product $\GK_c$ \\[3pt]
$Q$-enriched & $Q_{a+k} \cdot Q_{b+k}$ &
  Cross-correlations $\Phi_{a,b}$
\end{tabular}
\end{center}
The standard growth diagram produces the Borodin product formula
for $\GK_c$ (summing uniformly over all ALCDs), while the
$Q$-enriched growth diagram produces the cascade product
$\mathbf{Q}^*(w)\, \mathbf{Q}(z)$, whose coefficients are the
$Q$-kernels $\Phi_{a,b}(zw)$.
\end{remark}

\subsection{The commutator as obstruction}

The commutator $[\mathbf{Q}(z),\, \mathbf{Q}^*(w)]$
from~\eqref{eq:comm-Phi} measures the extent to which the
$Q$-enriched creation ($\mathbf{Q}$) and annihilation
($\mathbf{Q}^*$) fail to commute:
\begin{equation}\label{eq:comm-obstruction}
  [\mathbf{Q}(z),\, \mathbf{Q}^*(w)]
  = -\sum_{a,b \geq 0} P_a\, T_b\; z^a w^b\;
    \bigl(\Phi_{a,b}(zw) - Q_a Q_b\bigr).
\end{equation}
The obstruction vanishes (the operators commute) if and only if
$\Phi_{a,b}(t) = Q_a Q_b$ for all $a, b$, which happens precisely
when $Q_n = 0$ for all $n \geq 1$.  In the bijective framework,
the commutator terms with $k \geq 1$ correspond to face labels
that create ``cross-level interactions'' between different depths
of the cylindric partition.


%% ============================================================
\section{Special cases and examples}
\label{sec:special-cases}
%% ============================================================


\subsection{$d = 1$: recovery of the standard rule}

For $d = 1$ (profiles $(1,0,0)$ and $(0,1,0)$), $Q_n(q) = 1$
for all~$n$.  The $Q$-kernel becomes:
\begin{equation}\label{eq:Phi-d1}
  \Phi_{a,b}(t) = \sum_{k \geq 0} 1 \cdot 1 \cdot t^k
  = \frac{1}{1-t}.
\end{equation}
This is independent of $a$ and $b$: every face carries the same
generating function.  The $Q$-enriched local rule reduces to the
standard Burge local rule, confirming internal consistency.


\subsection{$d = 2$ (Rogers--Ramanujan)}

For $c = (1,1,0)$, $Q_n = q^{n^2}$ (the first Rogers--Ramanujan
profile).  The face weight for cascade parameter~$k$ is:
\begin{equation}\label{eq:face-weight-RR}
  \mathrm{wt}(k;\, a, b) = q^{(a+k)^2 + (b+k)^2}.
\end{equation}
The $Q$-kernel is:
\begin{equation}\label{eq:Phi-RR}
  \Phi_{a,b}(t)
  = \sum_{k \geq 0} q^{(a+k)^2 + (b+k)^2}\, t^k
  = q^{a^2 + b^2} \sum_{k \geq 0}
    q^{2k^2 + 2k(a+b)}\, t^k.
\end{equation}

\begin{example}[$\Phi_{0,0}$ and $\Phi_{1,0}$ for $d = 2$]
\label{ex:Phi-d2}
\begin{align}
  \Phi_{0,0}(t)
  &= \sum_{k \geq 0} q^{2k^2}\, t^k
  = 1 + q^2\, t + q^8\, t^2 + q^{18}\, t^3 + \cdots,
    \label{eq:Phi00-d2} \\
  \Phi_{1,0}(t)
  &= q\, \sum_{k \geq 0} q^{2k^2 + 2k}\, t^k
  = q + q^5\, t + q^{13}\, t^2 + \cdots.
    \label{eq:Phi10-d2}
\end{align}
The exponents $2k^2$ in $\Phi_{0,0}$ and $1 + 2k^2 + 2k
= (2k+1)^2/2 + 1/2$ in $\Phi_{1,0}$ reflect the
\emph{quadratic form} structure underlying the Rogers--Ramanujan
identities.
\end{example}

For the second profile $c = (2,0,0)$, $Q_n = q^{n(n+1)}$, and
the face weight is $q^{(a+k)(a+k+1) + (b+k)(b+k+1)}$.  The same
quadratic structure appears.

\subsection{$d = 4$, $c = (2,2,0)$}

Here $Q_0 = 1$, $Q_1 = q^2 + q^3 + q^4 + q^5$, and the
$Q$-kernel is no longer a monomial product.

\begin{example}[$\Phi_{1,0}$ and $\Phi_{1,1}$ for $c = (2,2,0)$]
\label{ex:Phi-d4}
\begin{align}
  \Phi_{1,0}(t) &= C_1(t)
  = Q_1 + Q_2\, Q_1\, t + Q_3\, Q_2\, t^2 + \cdots
    \notag \\
  &= (q^2 + q^3 + q^4 + q^5)
    + Q_2(q^2 + q^3 + q^4 + q^5)\, t + \cdots,
    \label{eq:Phi10-d4} \\[6pt]
  \Phi_{1,1}(t)
  &= Q_1^2 + Q_2^2\, t + Q_3^2\, t^2 + \cdots
    \notag \\
  &= (q^4 + 2q^5 + 3q^6 + 4q^7 + 3q^8 + 2q^9 + q^{10})
    + Q_2^2\, t + \cdots.
    \label{eq:Phi11-d4}
\end{align}
Note that $Q_1^2 = (q^2 + q^3 + q^4 + q^5)^2$ is itself a
polynomial with positive coefficients---this is a consequence of
$Q_1 \in \Z_{\geq 0}[q]$ (confirmed by computation for this
profile).
\end{example}


%% ============================================================
\section{Connection to the involution}
\label{sec:connection-involution}
%% ============================================================

\subsection{The peel/glue mechanism with $Q$-weights}

The generalised local rule provides a concrete mechanism for the
sign-reversing involution approach to proving positivity.  Recall
from~\cite{UNDERLYING_SET} that the ground set
$\mathcal{G}_{n,c}$ for the involution consists of pairs
$(S, \mathfrak{d})$ where $S \subseteq \{1, \ldots, n\}$ with
$|S| = j$ and $\mathfrak{d} \in \ALCD_{n-j}(c)$, with sign
$(-1)^j$ and weight $q^{\sigma(S) + |\mathfrak{d}|}$.

In the $Q$-enriched framework, the face labels of $\mathfrak{d}$
carry the additional weight $\prod_f Q_{a_f + \mathfrak{d}(f)}
\cdot Q_{b_f + \mathfrak{d}(f)}$.  The involution must pair
elements with opposite signs, and the unpaired elements (fixed
points with $j = 0$) should contribute exactly $Q_{n,c}(q)$.

\subsection{The $Q$-kernel as coupling constant}

The $Q$-kernel $\Phi_{a,b}(t)$ quantifies the coupling between
different depth levels.  The $k = 0$ term $Q_a Q_b$ represents
the ``uncoupled'' contribution (the product of independent
$Q$-weights), while the terms with $k \geq 1$ represent
\emph{cross-level interactions}.

From~\eqref{eq:comm-Phi}, the commutator isolates precisely these
cross-level terms:
\[
  \Phi_{a,b}(t) - Q_a Q_b
  = t\, \Phi_{a+1, b+1}(t)
  = \sum_{k \geq 1} Q_{a+k}\, Q_{b+k}\, t^k.
\]
In the language of the involution: the terms with $k \geq 1$ are
the contributions that the sign-reversing involution must cancel.
The commutator provides the algebraic bookkeeping for these
cancellations.

\subsection{Non-negativity and the conjecture}

\begin{proposition}\label{prop:Phi-nonneg}
If $Q_n(q) \in \Z_{\geq 0}[q]$ for all $n$ (Warnaar's
Conjecture~2.7), then for all $a, b, k$:
\[
  Q_{a+k}(q) \cdot Q_{b+k}(q) \in \Z_{\geq 0}[q],
\]
and consequently every coefficient of $\Phi_{a,b}(t)$ lies in
$\Z_{\geq 0}[q]$.  All face labels in the $Q$-enriched growth
diagram carry non-negative $q$-weights.
\end{proposition}

This is immediate from the closure of $\Z_{\geq 0}[q]$ under
multiplication, but it highlights an important structural point:
the positivity conjecture ensures that the $Q$-enriched growth
diagram has a combinatorial interpretation (all weights are
non-negative).  Conversely, if any $Q_n$ has a negative
coefficient, the corresponding face weights in the $Q$-enriched
diagram become signed, and the combinatorial interpretation
breaks down.

\subsection{Open questions}

\begin{enumerate}
  \item \textbf{Connecting the two frameworks.}
    Part~I defines $\hat{Q}_{c,\mathbf{n}}$ (an operator whose
    \emph{trace} gives $Q_{\mathbf{n},c}$), while Part~II defines
    $\mathbf{Q}(z)$ (an operator whose \emph{coefficients} are the
    $Q_{n,c}$).  Is there a natural relationship between
    $\hat{Q}_{c,\mathbf{n}}$ and the mode operators $Q_m P_m$?

  \item \textbf{Positive semi-definiteness.}
    Can the positivity of $Q_{n,c}(q)$ be deduced from
    positive semi-definiteness of $\hat{Q}^{\mathrm{sym}}$ or
    from positive-definiteness of
    $\mathbf{Q}^*(w)\mathbf{Q}(z)$?

  \item \textbf{Closed forms for $\Phi_{a,b}(t)$.}
    The $Q$-kernel $\Phi_{a,b}(t) = \sum Q_{a+k} Q_{b+k}\, t^k$
    involves shifted products of $Q$-polynomials.  For $d = 2$, it
    reduces to a theta-like sum~\eqref{eq:Phi-RR}.  Does
    $\Phi_{a,b}(t)$ have a closed-form expression in terms of the
    Borodin product or other known $q$-series for general
    profiles?

  \item \textbf{Matrix elements and the involution.}
    The matrix elements
    $\langle s_\lambda \mid P_a T_b \mid s_\mu \rangle$ count
    pairs of horizontal strips with common inner shape.  Can these
    be used to construct the sign-reversing involution, using the
    $Q$-kernel as the coupling constant?

  \item \textbf{The $(r,t)$-extension.}
    Both frameworks should extend to the Macdonald setting
    (replacing Schur by Macdonald polynomials, Pieri coefficients
    by $\varphi_{\lambda/\mu}(r,t)$ and
    $\psi_{\lambda/\mu}(r,t)$).
    Does the $Q$-kernel acquire additional structure in the
    $(r,t)$-deformed case?

  \item \textbf{Explicit $Q$-enriched bijections.}
    Can the $Q$-enriched growth diagram be used to construct
    explicit bijections between sets of cylindric partitions
    weighted by $Q$-polynomials?  In particular, does the
    bijection $\psi_c$ from~\eqref{eq:bijection-std} extend to a
    $Q$-weighted version that directly produces the
    cross-correlations $\Phi_{a,b}$?

  \item \textbf{The $Q$-kernel for $d \equiv 0 \pmod{3}$.}
    For $d$ divisible by~$3$, the polynomials $Q_n$ are defined
    using $(q^\ell;\, q^\ell)_n$ with $\ell = 3$.  The
    $Q$-kernel $\Phi_{a,b}(t)$ still makes sense formally, but the
    cyclic descent injection (see~\cite{FIRST_ATTEMPT}) fails.
    Does the $Q$-kernel reveal why $d \equiv 0$ is special?
\end{enumerate}


\begin{thebibliography}{99}
\bibitem{Borodin}
  A.~Borodin,
  \emph{Periodic Schur process and cylindric partitions},
  Duke Math.~J.\ \textbf{140} (2007), 391--468.

\bibitem{COMMUTATION}
  R.~Langer,
  \emph{From Commutation Relations to Growth Diagrams:
  A Summary of the Operator--Bijection Correspondence},
  companion note.

\bibitem{CONJECTURE}
  R.~Langer,
  \emph{The Positivity Conjecture for $Q_{n,c}(q)$},
  companion note.

\bibitem{FIRST_ATTEMPT}
  R.~Langer,
  \emph{Towards an Involution Proof of Positivity for $Q_{n,c}(q)$},
  companion note.

\bibitem{UNDERLYING_SET}
  R.~Langer,
  \emph{The Underlying Set: Cylindric Plane Partitions, ALCDs,
  and the Ground Set for the Involution},
  companion note.

\bibitem{langer-masters}
  R.~Langer,
  \emph{Symmetric Functions and Macdonald Polynomials},
  Master's thesis, University of Melbourne, 2008.
  arXiv:0907.3950.

\bibitem{langer-cpp1}
  R.~Langer,
  \emph{Enumeration of Cylindric Plane Partitions -- Part~I},
  2012. arXiv:1204.4583.

\bibitem{langer-thesis}
  R.~Langer,
  \emph{Cylindric Partitions and Bijective Proofs},
  PhD thesis, University of Melbourne.

\bibitem{Macdonald}
  I.~G.~Macdonald,
  \emph{Symmetric Functions and Hall Polynomials},
  2nd ed., Oxford University Press, 1995.

\bibitem{Welsh}
  T.~A.~Welsh,
  \emph{Unpublished notes on cylindric partition polynomiality}.
\end{thebibliography}

\end{document}
